\documentclass[10pt,letterpaper]{article}
\usepackage[margin=0.75in]{geometry}
\usepackage[T1]{fontenc}
\usepackage{lmodern,booktabs,amsmath}
\usepackage[hidelinks]{hyperref}
\setlength{\parindent}{0pt}
\setlength{\parskip}{5pt}
\title{\vspace{-1.8em}HUBO Degree Reduction for Pangenome-Guided Assembly\\
\large Concise Results Summary}
\author{Aarav Shandilya}
\date{}
\begin{document}
\maketitle
\vspace{-1.5em}
\section*{Purpose and scope}
The reported experiments reduce eighth-order higher-order unconstrained binary
optimisation (HUBO) models to maximum degree six (D6) or five (D5).
Three graph methods---Minigraph, Minigraph-Cactus, and PGGB---are compared at
12, 16, and 20 original binary variables, giving 18 configurations.
The results below are transcribed from the supplied presentation
(slides 2--7) \cite{slides}; they are not a new execution of the pipeline.
The background paper motivates the broader trade-off between compact HUBO
encodings and quantum circuit complexity \cite{cudby}.

\section*{Computational results}
Reduction decreases the largest number of variables interacting in a term,
but introduces auxiliary binary variables. The same variable counts are
reported for all three graph methods.
\begin{center}
\small
\begin{tabular}{rrrrrr}
\toprule
Original variables & Degree change & Auxiliaries & Total variables & Overhead & Space multiplier\\
\midrule
12 & $8\to6$ & 3 & 15 & 25.0\% & $8\times$\\
12 & $8\to5$ & 5 & 17 & 41.7\% & $32\times$\\
16 & $8\to6$ & 4 & 20 & 25.0\% & $16\times$\\
16 & $8\to5$ & 7 & 23 & 43.8\% & $128\times$\\
20 & $8\to6$ & 5 & 25 & 25.0\% & $32\times$\\
20 & $8\to5$ & 9 & 29 & 45.0\% & $512\times$\\
\bottomrule
\end{tabular}
\end{center}
\textbf{Scaling explanation.}
For $q$ original variables and $a$ auxiliaries, the number of possible binary
assignments increases from $2^q$ to $2^{q+a}$, a factor of $2^a$.
These multipliers are calculated from the reported variable counts;
they are not measured runtime increases. Full statevector storage also
scales as $2^{q+a}$, although other solvers need not store every state.
At 20 original variables, D5 has $2^{29-25}=16$ times as many assignments
as D6. D6 therefore has lower variable overhead throughout the tested sizes.

\section*{Main biological findings}
Fifteen of the 18 configurations have no reported change in identity,
mismatches, insertions, or deletions. The three exceptions improve identity:
\begin{itemize}
\setlength{\itemsep}{0pt}
\item \textbf{Minigraph, 16 variables, D5:} identity increases by 4.64 percentage
points, with 33 fewer mismatches and two fewer deletions, but one additional insertion.
\item \textbf{PGGB and Minigraph-Cactus, 16 variables, D6:} each improves identity
by 0.25 percentage points, with two fewer mismatches and no insertion/deletion change.
\end{itemize}
No displayed configuration loses identity, but neither degree consistently
improves it. These results concern changes in the reported sequence metrics;
unchanged metrics alone do not prove that the decoded sequence or graph walk is identical.

\textbf{Simple explanation.}
Think of a genome as an essay assembled from linked sentences. Reduction
simplifies the rules connecting the sentences, but adds extra bookkeeping
choices. Most reported essays retain the same accuracy; three improve.
More bookkeeping can make searching for the best essay harder.

\newpage
\section*{Complete biological comparison}
All changes are reduced minus unreduced. Identity is in percentage points
(pp); negative mismatch, insertion, or deletion changes mean fewer reported
errors. The presentation does not specify whether insertion/deletion counts
represent events or affected bases, so that unit is left unspecified here.
Original denotes the original binary-variable count, labelled Qubits
in the slides; it is not evidence of a hardware run.
\begin{center}
\small
\setlength{\tabcolsep}{5pt}
\begin{tabular}{llrrrrr}
\toprule
Method & Target & Original & $\Delta$ Identity (pp) & $\Delta$ Mismatch & $\Delta$ Insertion & $\Delta$ Deletion\\
\midrule
Minigraph & D6 & 12 & 0.00 & 0 & 0 & 0\\
Minigraph & D5 & 12 & 0.00 & 0 & 0 & 0\\
Minigraph & D6 & 16 & 0.00 & 0 & 0 & 0\\
Minigraph & D5 & 16 & +4.64 & $-33$ & +1 & $-2$\\
Minigraph & D6 & 20 & 0.00 & 0 & 0 & 0\\
Minigraph & D5 & 20 & 0.00 & 0 & 0 & 0\\
\midrule
Minigraph-Cactus & D6 & 12 & 0.00 & 0 & 0 & 0\\
Minigraph-Cactus & D5 & 12 & 0.00 & 0 & 0 & 0\\
Minigraph-Cactus & D6 & 16 & +0.25 & $-2$ & 0 & 0\\
Minigraph-Cactus & D5 & 16 & 0.00 & 0 & 0 & 0\\
Minigraph-Cactus & D6 & 20 & 0.00 & 0 & 0 & 0\\
Minigraph-Cactus & D5 & 20 & 0.00 & 0 & 0 & 0\\
\midrule
PGGB & D6 & 12 & 0.00 & 0 & 0 & 0\\
PGGB & D5 & 12 & 0.00 & 0 & 0 & 0\\
PGGB & D6 & 16 & +0.25 & $-2$ & 0 & 0\\
PGGB & D5 & 16 & 0.00 & 0 & 0 & 0\\
PGGB & D6 & 20 & 0.00 & 0 & 0 & 0\\
PGGB & D5 & 20 & 0.00 & 0 & 0 & 0\\
\bottomrule
\end{tabular}
\end{center}

\section*{Interpretation and next steps}
The displayed results support preservation of the measured assembly accuracy
in most configurations, with occasional improvement. D6 uses fewer
auxiliaries; D5 offers a lower maximum interaction degree but a larger search
space. These tables alone do not establish which degree is faster, requires
shallower circuits, or performs better on quantum hardware. Improvement
in sampled solutions also does not, by itself, prove mathematical equivalence
of the reduced and original objectives.

The supplied comparison lacks absolute identity values, coverage, assembly
break counts, timings, circuit statistics, and repeated-run uncertainty.
Accordingly, it cannot establish complete genome reconstruction or a
statistically reliable advantage. The next useful checks are to validate
legal graph walks and auxiliary consistency, compare decoded sequences
against ground truth, and measure runtime and circuit resources across
repeated runs before assessing IonQ hardware performance.

\begin{thebibliography}{9}
\small
\bibitem{slides}
A.~Shandilya, \emph{Pangenome-Guided Genome Assembly: Sixth-Order HUBO
Reduction and Biological Validation}, supplied presentation,
\emph{Copy of HUBO Reduction - IonQ Presentation(1).pdf}, slides 2--8.
\bibitem{cudby}
J.~Cudby and S.~Strelchuk, \emph{Nonvariational quantum optimisation
approaches to pangenome-guided sequence assembly},
arXiv:2604.06106v1, 2026. Supplied PDF: \emph{NonVariational .pdf}.
\url{https://arxiv.org/abs/2604.06106}
\end{thebibliography}
\end{document}
